%% Abstract (configurado para língua inglesa)
\begin{resumo}[Abstract]			% Título do Resumo (Abstract = Resumo em inglês)
\begin{otherlanguage*}{english}		% Língua do texto
Currently, errors in the calculation of criminal enforcement have contributed to the undue
retention of convicts in the Brazilian prison system, aggravating a scenario already marked
by an increase in the prison population. Among the factors involved, one of the most
recurring is errors in the Criminal Enforcement Process (PEC). The PEC is the
set of procedures that accompany the serving of sentences, including calculations,
the granting of benefits, and the recording of work and/or study activities for remission,
in accordance with the Criminal Enforcement Law (LEP). When performed manually, this calculation
is susceptible to errors that go unnoticed, which can result in prison terms
longer than those legally prescribed or in undue early releases,
affecting the management of prison places and public safety. The correct enforcement of the sentence and
the effective calculation, considering the remissions obtained by the convict, contribute
to the individual's progress in the resocialization process, reducing recidivism
and contributing to the safety of society. Thus, the objective of this project is to
automate the calculation of criminal enforcement, reducing human error and ensuring greater
accuracy and agility in the enforcement of sentences.

Translated with DeepL.com (free version)

\vspace{\onelineskip}
\noindent
\textbf{Keywords}: criminal enforcement; prison system; automation; automatic calculation;
Criminal Enforcement Law; PEC.
\end{otherlanguage*}
\end{resumo}

%% Exemplo de resumo em francês
%\begin{resumo}[Résumé]
% \begin{otherlanguage*}{french}
%    Il s'agit d'un résumé en français.
% 
%   \textbf{Mots-clés}: latex. abntex. publication de textes.
% \end{otherlanguage*}
%\end{resumo}

%% Exemplo de resumo em Espanhol
%\begin{resumo}[Resumen]
% \begin{otherlanguage*}{spanish}
%   Este es el resumen en español.
%  
%   \textbf{Palabras clave}: latex. abntex. publicación de textos.
% \end{otherlanguage*}
%\end{resumo}
% ---